% ============================================================
%  MTH636 -- Game Theory | Lecture 7
%  Nash Equilibrium & Best Reply Correspondence
% ============================================================
\documentclass[12pt, a4paper]{article}

% ============================================================
%  MTH636 -- Game Theory
%  Notes Preamble — Color-Coded, TikZ-Enhanced
% ============================================================

% --- Core packages ---
\usepackage[T1]{fontenc}
\usepackage[utf8]{inputenc}
\usepackage{lmodern}
\usepackage[margin=2.5cm, top=2.8cm, bottom=2.8cm]{geometry}
\usepackage{microtype}

% --- Math ---
\usepackage{amsmath, amssymb, amsthm, mathtools}
\usepackage{bm}        % bold math
\usepackage{cancel}    % strike-through in math

% --- Color & Boxes ---
\usepackage[dvipsnames, svgnames, x11names, table]{xcolor}
\usepackage[most]{tcolorbox}
\tcbuselibrary{breakable, skins, theorems}

% --- TikZ & Diagrams ---
\usepackage{tikz}
\usepackage{tikz-qtree}
\usetikzlibrary{arrows.meta, positioning, shapes.geometric, fit, backgrounds,
                decorations.pathreplacing, calc, matrix, trees}
\usepackage{pgfplots}
\pgfplotsset{compat=1.18}

% --- Tables ---
\usepackage{booktabs}
\usepackage{multirow}
\usepackage{array}
\usepackage{tabularx}
\usepackage{makecell}

% --- Typography & Layout ---
\usepackage{enumitem}
\usepackage{parskip}
\usepackage{fancyhdr}
\usepackage{titlesec}
\usepackage{titletoc}
\usepackage{hyperref}
\usepackage{graphicx}
\usepackage{float}
\usepackage{caption}

% --- Color Palette ---
\definecolor{defblue}{HTML}{1A5276}      % Definition header
\definecolor{defbluebg}{HTML}{D6EAF8}    % Definition background
\definecolor{thmgreen}{HTML}{1E8449}     % Theorem header
\definecolor{thmbg}{HTML}{D5F5E3}        % Theorem background
\definecolor{exorange}{HTML}{B7770D}     % Example header
\definecolor{exbg}{HTML}{FEF9E7}         % Example background
\definecolor{warnred}{HTML}{922B21}      % Warning/Key Insight header
\definecolor{warnbg}{HTML}{FADBD8}       % Warning background
\definecolor{noticegray}{HTML}{2C3E50}   % Note header
\definecolor{noticebg}{HTML}{EBF5FB}     % Note background
\definecolor{taketawayyellow}{HTML}{7D6608} % Key Takeaway header
\definecolor{takeawaybg}{HTML}{FDFEFE}   % Takeaway background
\definecolor{accentpurple}{HTML}{6C3483} % Accent / proof
\definecolor{lightgray}{HTML}{F2F3F4}    % shading

% ============================================================
%  TCOLORBOX ENVIRONMENTS
% ============================================================

% -- Definition (Blue) --
\newtcolorbox{defbox}[2][]{
  enhanced, breakable,
  colback=defbluebg, colframe=defblue, coltitle=white,
  fonttitle=\bfseries\sffamily,
  title={Definition: #2},
  attach boxed title to top left={yshift=-2mm, xshift=4mm},
  boxed title style={colback=defblue, rounded corners},
  arc=4pt, boxrule=1.2pt,
  left=6pt, right=6pt, top=8pt, bottom=6pt,
  #1
}

% -- Theorem / Proposition (Green) --
\newtcolorbox{thmbox}[2][]{
  enhanced, breakable,
  colback=thmbg, colframe=thmgreen, coltitle=white,
  fonttitle=\bfseries\sffamily,
  title={Theorem: #2},
  attach boxed title to top left={yshift=-2mm, xshift=4mm},
  boxed title style={colback=thmgreen, rounded corners},
  arc=4pt, boxrule=1.2pt,
  left=6pt, right=6pt, top=8pt, bottom=6pt,
  #1
}

% -- Proposition (Green variant) --
\newtcolorbox{propbox}[2][]{
  enhanced, breakable,
  colback=thmbg, colframe=thmgreen!70!black, coltitle=white,
  fonttitle=\bfseries\sffamily,
  title={Proposition: #2},
  attach boxed title to top left={yshift=-2mm, xshift=4mm},
  boxed title style={colback=thmgreen!70!black, rounded corners},
  arc=4pt, boxrule=1.2pt,
  left=6pt, right=6pt, top=8pt, bottom=6pt,
  #1
}

% -- Example (Orange/Yellow) --
\newtcolorbox{exbox}[2][]{
  enhanced, breakable,
  colback=exbg, colframe=exorange, coltitle=white,
  fonttitle=\bfseries\sffamily,
  title={Example: #2},
  attach boxed title to top left={yshift=-2mm, xshift=4mm},
  boxed title style={colback=exorange, rounded corners},
  arc=4pt, boxrule=1.2pt,
  left=6pt, right=6pt, top=8pt, bottom=6pt,
  #1
}

% -- Key Insight / Warning (Red) --
\newtcolorbox{warnbox}[2][]{
  enhanced, breakable,
  colback=warnbg, colframe=warnred, coltitle=white,
  fonttitle=\bfseries\sffamily,
  title={Key Insight: #2},
  attach boxed title to top left={yshift=-2mm, xshift=4mm},
  boxed title style={colback=warnred, rounded corners},
  arc=4pt, boxrule=1.2pt,
  left=6pt, right=6pt, top=8pt, bottom=6pt,
  #1
}

% -- Note (Gray-Blue) --
\newtcolorbox{notebox}[1][]{
  enhanced, breakable,
  colback=noticebg, colframe=noticegray!60, coltitle=white,
  fonttitle=\bfseries\sffamily,
  title={Note},
  attach boxed title to top left={yshift=-2mm, xshift=4mm},
  boxed title style={colback=noticegray!60, rounded corners},
  arc=4pt, boxrule=0.8pt,
  left=6pt, right=6pt, top=8pt, bottom=6pt,
  #1
}

% -- Key Takeaway --
\newtcolorbox{takebox}[1][]{
  enhanced, breakable,
  colback=takeawaybg, colframe=taketawayyellow, coltitle=white,
  fonttitle=\bfseries\sffamily,
  title={Key Takeaway},
  attach boxed title to top left={yshift=-2mm, xshift=4mm},
  boxed title style={colback=taketawayyellow, rounded corners},
  arc=4pt, boxrule=1.2pt,
  left=6pt, right=6pt, top=8pt, bottom=6pt,
  #1
}

% -- Intuition --
\newtcolorbox{intubox}[1][]{
  enhanced, breakable,
  colback=accentpurple!8, colframe=accentpurple!60, coltitle=white,
  fonttitle=\bfseries\sffamily,
  title={Intuition},
  attach boxed title to top left={yshift=-2mm, xshift=4mm},
  boxed title style={colback=accentpurple!80, rounded corners},
  arc=4pt, boxrule=0.8pt,
  left=6pt, right=6pt, top=8pt, bottom=6pt,
  #1
}

% ============================================================
%  PAYOFF MATRIX HELPER
%  Usage: inside tabular, use \payoff{a}{b} for (a,b)
% ============================================================
\newcommand{\payoff}[2]{$(#1,\;#2)$}
\newcommand{\nash}[2]{\underline{$#1$},\underline{$#2$}}  % underline NE payoffs
\newcommand{\Nset}{\mathbb{N}}
\newcommand{\Rset}{\mathbb{R}}
\newcommand{\Sset}{\mathcal{S}}
\newcommand{\Ui}{u_i}

% ============================================================
%  SECTION FORMATTING
% ============================================================
\titleformat{\section}
  {\large\bfseries\sffamily\color{defblue}}
  {\thesection}{1em}{}[\titlerule]

\titleformat{\subsection}
  {\normalsize\bfseries\sffamily\color{thmgreen!80!black}}
  {\thesubsection}{1em}{}

\titleformat{\subsubsection}
  {\normalsize\itshape\sffamily\color{exorange!80!black}}
  {\thesubsubsection}{1em}{}

% ============================================================
%  HEADER / FOOTER
% ============================================================
\pagestyle{fancy}
\fancyhf{}
\fancyhead[L]{\sffamily\small\color{defblue} MTH636 -- Game Theory}
\fancyhead[R]{\sffamily\small\color{defblue} \leftmark}
\fancyfoot[C]{\sffamily\small\thepage}
\renewcommand{\headrulewidth}{0.4pt}

% ============================================================
%  HYPERREF SETUP
% ============================================================
\hypersetup{
  colorlinks=true,
  linkcolor=defblue,
  urlcolor=thmgreen,
  citecolor=accentpurple,
  bookmarksnumbered=true,
  pdftitle={MTH636 Game Theory Notes},
  pdfauthor={MTH636 Student Notes}
}

% ============================================================
%  MISC HELPERS
% ============================================================
\newcommand{\BR}{\text{BR}}        % Best Response
\newcommand{\NE}{\text{NE}}        % Nash Equilibrium
\newcommand{\SPNE}{\text{SPNE}}    % Subgame Perfect NE
\newcommand{\Dom}{\text{dom}}      % Dominance
\newcommand{\mixs}[1]{\sigma_{#1}} % Mixed strategy
\newcommand{\pures}[1]{s_{#1}}     % Pure strategy


% ============================================================
\begin{document}

% Title
\begin{center}
  {\LARGE\bfseries\sffamily\color{defblue} MTH636 -- Game Theory}\\[6pt]
  {\Large\sffamily Lecture 7: Nash Equilibrium}\\[4pt]
\end{center}
\vspace{0.5cm}
\hrule height 1.5pt
\vspace{0.3cm}

% ============================================================
\section{Formal Definition of Nash Equilibrium}

\begin{defbox}{Nash Equilibrium}
  For a game $G = \langle N, (S_i)_{i \in N}, (u_i)_{i \in N} \rangle$, a strategy profile $s^*$ is a \textbf{Nash Equilibrium (NE)} if for all $i \in N$ and for all $s_i \in S_i$:
  \[
    u_i(s^*) \geq u_i(s_i, s^*_{-i})
  \]
  In words: No player $i$ can strictly increase their payoff by unilaterally changing their strategy from $s_i^*$ to some other $s_i$, given that all other players play their part of $s^*_{-i}$.
\end{defbox}

\begin{exbox}{Example: Battle of the Sexes (Variants)}
  Consider the following game matrix:
  \begin{center}
  \renewcommand{\arraystretch}{1.3}
  \begin{tabular}{c|cc}
    $P_1 \setminus P_2$ & F & D \\
    \hline
    F & \payoff{2}{1} & \payoff{0}{0} \\
    D & \payoff{0}{0} & \payoff{1}{2} \\
  \end{tabular}
  \end{center}
  Here, both $(F, F)$ and $(D, D)$ are Nash Equilibria. If they are at $(F, F)$, neither wants to deviate to $D$ (yielding 0).
\end{exbox}

% ============================================================
\section{Best Reply Correspondence}

To find Nash Equilibria, it is often useful to look at players' best responses.

\begin{defbox}{Best Reply Correspondence}
  The best reply correspondence for player $i$ is a function $b_i: S_{-i} \to 2^{S_i}$ (mapping to the power set of $S_i$), defined as:
  \[
    b_i(s_{-i}) = \arg\max_{t_i \in S_i} u_i(t_i, s_{-i})
  \]
  This gives the set of all strategies $t_i$ that maximize player $i$'s payoff given the opponents' profile $s_{-i}$.
\end{defbox}

\begin{thmbox}{Nash Equilibrium in terms of Best Replies}
  A strategy profile $s^*$ is a Nash Equilibrium if and only if:
  \[
    s_i^* \in b_i(s^*_{-i}) \quad \forall i \in N
  \]
\end{thmbox}

% ============================================================
\section{Finding Nash Equilibria in Specific Games}

\subsection{Beauty Contest Game}
Recall the Beauty Contest Game: $N = \{1, 2, 3\}$, $S_i = [0, 100]$.
Target $T(x_1, x_2, x_3) = \frac{2}{3} \left( \frac{x_1 + x_2 + x_3}{3} \right)$.
Players who are closest to $T$ win. 

Is the profile $(0, 0, 0)$ a Nash Equilibrium? Yes, because if everyone plays 0, the target is 0, and deviating yields a strictly positive number which is further from 0.

What about profiles of the form $(a, a, a)$ where $a \in (0, 100]$? These are \textbf{not} Nash Equilibria, because any player could deviate to a slightly lower number to get closer to $\frac{2}{3}$ of the average and win outright.

\subsection{Second-Price Auction}
Recall that bidding true valuation weakly dominates all other strategies. Therefore:
\begin{itemize}
    \item The profile $(v_1, v_2, \dots, v_n)$ is a Nash Equilibrium.
    \item What about $(v_1, 0, 0, \dots, 0)$? If $v_1$ is the highest valuation, Player 1 wins and pays 0. No other player can profitably deviate (bidding higher than $v_1$ means they win but pay more than their valuation, giving a negative payoff). So this can also be an NE under certain conditions.
\end{itemize}
\textbf{Exercise:} Find all Nash Equilibria of the Second-Price Auction game where $v_1 = 10, v_2 = 11$.

% ============================================================
\section{Continuous Action Spaces (Cournot-like Game)}

Consider a 2-player game where $N = \{F_1, F_2\}$ and $S_i = [0, 1]$.
Payoffs are given by:
\[
  u_i(a_1, a_2) = (2 - a_1 - a_2)a_i - c_i a_i
\]
where $c_1, c_2$ are cost parameters.

To find the Nash Equilibrium $(a_1^*, a_2^*)$, we use the Best Reply functions.
Player 1 maximizes $u_1(a_1, \bar{a}_2)$ with respect to $a_1$:
\[
  \frac{\partial u_1}{\partial a_1} = 2 - 2a_1 - \bar{a}_2 - c_1 = 0 \implies a_1 = \frac{2 - c_1 - \bar{a}_2}{2}
\]
Thus, the best reply function is $b_1(\bar{a}_2) = \frac{2 - c_1 - \bar{a}_2}{2}$.
Similarly, for Player 2: $b_2(\bar{a}_1) = \frac{2 - c_2 - \bar{a}_1}{2}$.

At the Nash Equilibrium, both must be playing best responses to each other:
\[
  a_1^* = \frac{2 - c_1 - a_2^*}{2} \quad \text{and} \quad a_2^* = \frac{2 - c_2 - a_1^*}{2}
\]
Substitute $a_2^*$ into the equation for $a_1^*$:
\begin{align*}
  a_1^* &= \frac{2 - c_1 - \left( \frac{2 - c_2 - a_1^*}{2} \right)}{2} \\
  2a_1^* &= 2 - c_1 - 1 + \frac{c_2}{2} + \frac{a_1^*}{2} \\
  4a_1^* &= 2 - 2c_1 + c_2 + a_1^* \\
  3a_1^* &= 2 - 2c_1 + c_2 \implies a_1^* = \frac{2 - 2c_1 + c_2}{3}
\end{align*}
By symmetry, $a_2^* = \frac{2 - 2c_2 + c_1}{3}$.

The equilibrium payoffs are:
\[
  u_1(a_1^*, a_2^*) = (a_1^*)^2 \quad \text{and} \quad u_2(a_1^*, a_2^*) = (a_2^*)^2
\]

\end{document}